\section{Introducción: en 2025 el protagonismo pasa de las empresas de modelos a las aplicaciones y los Agent}

Esta sección sitúa el episodio. Zhang Xiaojun (张小珺) explica al inicio que en 2024 el relato de la IA todavía giraba en torno a las empresas de modelos grandes, y que en 2025 los nuevos protagonistas son las empresas de aplicaciones y los Agent. Ming Chaoping (明超平) es un emprendedor de aplicaciones de IA nacido en 1995; trabajó en producto sucesivamente en OnePlus, ByteDance y Moonshot, es CEO por primera vez y el primer producto que lanzó se llama YouWare. El valor de este episodio no está solo en escuchar la historia personal de un founder joven, sino en ver cómo un fundador con perfil de producto enlaza teléfonos móviles, video corto, Kimi, coding, Agent y la oportunidad del OS en un mismo conjunto de juicios.

\begin{importantbox}{Tesis central del episodio}
La oportunidad de YouWare no se reduce a «hacer un producto de Agent»: consiste en redefinir el contenedor de creación y el entorno de ejecución de la era de la IA. El juicio de Ming Chaoping es que los modelos base seguirán mejorando, pero las empresas de aplicaciones deben adaptar modelos más inteligentes a necesidades de producción concretas mediante el entorno, la experiencia, la interacción y los efectos de red.
\end{importantbox}

\begin{knowledgebox}{Nota sobre la estrategia visual}
Este video es una entrevista con encuadre fijo, sin slides, pizarrón ni demostración de producto. El texto usa la portada solo para identificar la fuente; todas las imágenes del cuerpo son diagramas conceptuales de elaboración propia, que explican las tres etapas como gerente de producto, experiencia y datos, la bitter lesson, el consumo de token, el contenedor/entorno, el Vibe Coder, AgentRank, el OS Agent y la confrontación con el Ego.
\end{knowledgebox}

\begin{figure}[H]
\centering
\includegraphics[width=0.96\textwidth]{product-manager-three-stations.png}
\caption{Tres etapas como gerente de producto: OnePlus, ByteDance y Moonshot forman la base de su juicio de producto. Diagrama conceptual de elaboración propia, organizado a partir de la conversación de 00:03:16--01:05:05.}
\end{figure}

\begin{knowledgebox}{Lectura de la figura: cada una de las tres experiencias entrena una capacidad de producto distinta}
OnePlus entrena la experiencia y la sensibilidad estética; ByteDance, los datos, el crecimiento y la iteración de alta frecuencia; Moonshot, los modelos, los productos de IA y la incertidumbre de la frontera. YouWare no surge de la nada: es la recombinación de estas tres capacidades en la ventana de oportunidad de las aplicaciones de Agent.
\end{knowledgebox}

\subsection{Por qué esta no es una entrevista de emprendimiento común}

Esta sección desplaza primero la atención de la trayectoria personal al método de producto. La historia de Ming Chaoping tiene dos líneas: una es cómo una persona forma su juicio de producto; la otra, cómo una empresa de aplicaciones de IA encuentra su lugar junto a los gigantes de los modelos. La primera incluye el entrenamiento en debate, la competencia de vehículos inteligentes, el entrenamiento en experiencia de OnePlus y el entrenamiento en datos de ByteDance; la segunda incluye la bitter lesson, el consumo de token, el entorno como contenedor, AgentRank y el OS Agent.

\begin{knowledgebox}{Nota de clase: leer la historia personal como metodología}
La lectura útil de este episodio no es «el founder joven tiene mucha personalidad», sino «qué experiencias lo llevaron a qué juicios de producto». El debate corresponde a la perspectiva de un tercero; la competencia de vehículos inteligentes, al reward y a la retroalimentación práctica; OnePlus, a la experiencia; ByteDance, a los datos; Moonshot, al sentido de los límites en la era de los modelos.
\end{knowledgebox}

\subsection{Ruta de lectura}

Esta sección propone cómo leer el episodio. La primera mitad de la entrevista trata del debate, OnePlus, ByteDance y Kimi, y es fácil malinterpretarla como un currículum personal; la segunda mitad trata de los token, el contenedor, el Vibe Coder, AgentRank y el OS Agent, que son el núcleo del juicio de producto. La lectura correcta es: la primera mitad explica por qué hace producto de esta manera; la segunda, por qué las empresas de aplicaciones de IA todavía tienen oportunidad.

\noindent\begin{tabularx}{\textwidth}{p{0.18\textwidth}p{0.34\textwidth}X}
\toprule
Pregunta de lectura & Material de la entrevista & Juicio que se debe formar \\
\midrule
De dónde viene la intuición de producto & Debate, OnePlus, ByteDance, Kimi & La moldean en conjunto la perspectiva de un tercero, la sensibilidad a la experiencia, la infraestructura de datos y la ventana de los modelos. \\
Cómo evaluar las aplicaciones de IA & bitter lesson, consumo de token, per token valuation & No pulir adornos; reconstruir el producto en torno a la capacidad del modelo y a la eficiencia con que se consume la inteligencia. \\
De dónde surge el ecosistema de Agent & AgentRank, efectos de red, documento de identidad, confianza & Los Agent no son herramientas aisladas: forman redes de invocación, clasificación y confianza. \\
Cómo construir los mecanismos del CEO & Valor emocional, andar como sobre hielo delgado, confrontar el Ego, mentalidad de jugador de ajedrez & Emprender en aplicaciones exige que el fundador se recalibre continuamente a sí mismo y a su organización. \\
\bottomrule
\end{tabularx}

\subsection{Resumen de la sección}

La línea principal de EP101 es que el emprendimiento en aplicaciones de IA pasa del culto a los modelos a la construcción de entornos. El caso de YouWare muestra que las empresas de aplicaciones no pueden limitarse a esperar a que los modelos se vuelvan más potentes: deben decidir en qué contenedor se coloca la inteligencia, cómo se consume y cómo genera efectos de red.
